BioMoAX organizes morphology, expression and genetic perturbation data within a unified retrieval framework, allowing the same encoders and drug reference library to support functional comparisons across modalities, experimental sources, cellular backgrounds and perturbation types. In A549, methods ranked differently in image-to-image and cross-modal MOA retrieval; BioMoAX showed its clearest advantage in functional matching between morphology and expression. Compound pairing provides identity correspondence between measurements, while shared MOAs and gene--drug associations supervise functional relationships across perturbations. This design allows jointly trained representations to retain drug identity and connect mechanism-related perturbations, providing a common functional reference for multiple query tasks.

This reference-based approach is useful when measurements are incomplete across experiments. When only imaging is available for a cellular background, cross-cell-line image-to-expression retrieval can identify mechanism-related responses in existing transcriptomic collections, informing pathway hypotheses and follow-up validation. Genetic queries extend this use to gene-function studies: existing drug responses provide a common reference for comparing gene perturbations and connecting genetic-screen findings to specific chemical interventions. Independent literature supporting the two co-retrieval cases shows how this reference can complement relationships absent from current gene--drug annotations. Existing morphology, expression and drug measurements can thus remain useful in subsequent experiments, providing a comparable functional context for further queries.

Current evaluation is limited by cellular backgrounds, genetic perturbation numbers, and annotation coverage. Crop averaging loses cellular heterogeneity, while dose and exposure time remain unmodeled. Broader gene and compound coverage and representations retaining cellular variation should clarify how functional relationships vary across doses, times, and backgrounds. Genetic and chemical perturbation experiments testing shared downstream responses would clarify the biological basis and scope of these retrieval relationships.
